\documentclass[11pt,a4paper]{article}

% ---- Packages ----
\usepackage[utf8]{inputenc}
\usepackage[T1]{fontenc}
\usepackage{amsmath,amssymb,amsthm}
\usepackage{mathtools}
\usepackage[margin=1in,headheight=14pt]{geometry}
\usepackage{hyperref}
\usepackage{enumitem}
\usepackage{xcolor}
\usepackage{tcolorbox}
\usepackage{fancyhdr}
\usepackage{booktabs}

% ---- Hyperref ----
\hypersetup{
    colorlinks=true,
    linkcolor=red,
    citecolor=blue,
    urlcolor=blue
}

% ---- Header ----
\pagestyle{fancy}
\fancyhf{}
\fancyhead[L]{\textit{Lesson 11: Static and Extensive-Form Games}}
\fancyhead[R]{\thepage}
\renewcommand{\headrulewidth}{0.4pt}

% ---- Theorem Environments ----
\theoremstyle{definition}
\newtheorem{definition}{Definition}[section]
\newtheorem{example}{Example}[section]
\newtheorem{exercise}{Exercise}[section]

\theoremstyle{plain}
\newtheorem{theorem}{Theorem}[section]
\newtheorem{lemma}[theorem]{Lemma}
\newtheorem{proposition}[theorem]{Proposition}
\newtheorem{corollary}[theorem]{Corollary}

\theoremstyle{remark}
\newtheorem{remark}{Remark}[section]

% ---- Colored Boxes ----
\tcbuselibrary{theorems,skins,breakable}

\newtcolorbox{keyresult}[1][]{
    colback=green!5!white,
    colframe=green!50!black,
    fonttitle=\bfseries,
    title=#1,
    breakable
}

\newtcolorbox{rlconnection}[1][]{
    colback=blue!5!white,
    colframe=blue!50!black,
    fonttitle=\bfseries,
    title=#1,
    breakable
}

\newtcolorbox{warning}[1][]{
    colback=red!5!white,
    colframe=red!50!black,
    fonttitle=\bfseries,
    title=#1,
    breakable
}

% ---- Operators ----
\newcommand{\R}{\mathbb{R}}
\newcommand{\N}{\mathbb{N}}
\newcommand{\Q}{\mathbb{Q}}
\newcommand{\Z}{\mathbb{Z}}
\newcommand{\C}{\mathbb{C}}
\newcommand{\Prob}{\mathbb{P}}
\newcommand{\E}{\mathbb{E}}
\newcommand{\indicator}{\mathbf{1}}
\DeclareMathOperator{\supp}{supp}
\DeclareMathOperator{\conv}{conv}
\DeclareMathOperator{\argmax}{argmax}
\DeclareMathOperator{\argmin}{argmin}
\DeclareMathOperator{\BR}{BR}

% ---- Title ----
\title{Lesson 34: Mixed Strategies, Nash Existence, and the Minimax Theorem}
\author{Curriculum: A Rigorous Learning Path to Multi-Agent Deep RL}
\date{}

\begin{document}
\maketitle
\tableofcontents
\newpage

\setcounter{section}{3}

\section{Mixed Strategies and Mixed Extensions}
\label{sec:mixed}

\subsection{Definitions}
\label{subsec:mixed_def}

\begin{definition}[Mixed Strategy]
\label{def:mixed_strategy}
Let $\Gamma$ be a finite game with strategy set $S_i = \{s_i^1, \ldots, s_i^{m_i}\}$
for player~$i$. A \emph{mixed strategy} for player~$i$ is a probability distribution
$\sigma_i \in \Delta(S_i)$, where
\[
    \Delta(S_i) = \bigg\{\sigma_i \in \R^{m_i} : \sigma_i(s_i^j) \geq 0
    \text{ for all } j, \;\; \sum_{j=1}^{m_i} \sigma_i(s_i^j) = 1\bigg\}
\]
is the \emph{$(m_i-1)$-simplex} over $S_i$. A \emph{mixed strategy profile} is
$\sigma = (\sigma_1, \ldots, \sigma_n) \in \Delta := \prod_{i \in N} \Delta(S_i)$.
\end{definition}

\begin{definition}[Support]
\label{def:support}
The \emph{support} of a mixed strategy $\sigma_i$ is
$\supp(\sigma_i) = \{s_i \in S_i : \sigma_i(s_i) > 0\}$.
A pure strategy $s_i$ is identified with the degenerate distribution placing
probability~1 on $s_i$.
\end{definition}

\begin{definition}[Mixed Extension]
\label{def:mixed_extension}
The \emph{mixed extension} of a finite game $\Gamma = (N, (S_i), (u_i))$ is the game
$\bar{\Gamma} = (N, (\Delta(S_i)), (\bar{u}_i))$ where the expected payoff under
independent mixing is
\[
    \bar{u}_i(\sigma) = \sum_{s \in S} \bigg(\prod_{j \in N} \sigma_j(s_j)\bigg) u_i(s)
    = \E_{s \sim \sigma}[u_i(s)].
\]
We drop the bar and write $u_i(\sigma)$ when the context is clear.
\end{definition}

\begin{remark}[Multilinearity]
\label{rem:multilinear}
The expected payoff $u_i(\sigma)$ is \emph{multilinear}: it is linear in each
$\sigma_i$ when the other players' strategies are held fixed. Explicitly,
\[
    u_i(\sigma_i, \sigma_{-i}) = \sum_{s_i \in S_i} \sigma_i(s_i)
    \underbrace{\sum_{s_{-i} \in S_{-i}}
    \bigg(\prod_{j \neq i} \sigma_j(s_j)\bigg) u_i(s_i, s_{-i})}_{=:\, u_i(s_i, \sigma_{-i})}.
\]
The quantity $u_i(s_i, \sigma_{-i})$ is the expected payoff to player~$i$ from playing
the pure strategy $s_i$ when opponents play $\sigma_{-i}$.
\end{remark}

\subsection{Mixed Nash Equilibrium and the Indifference Principle}
\label{subsec:mixed_ne}

\begin{definition}[Mixed Strategy Nash Equilibrium]
\label{def:mixed_ne}
A mixed strategy profile $\sigma^*$ is a \emph{Nash equilibrium} if
\[
    u_i(\sigma_i^*, \sigma_{-i}^*) \geq u_i(\sigma_i, \sigma_{-i}^*)
    \quad \text{for all } \sigma_i \in \Delta(S_i), \text{ for all } i \in N.
\]
\end{definition}

\begin{theorem}[Indifference Principle]
\label{thm:indifference}
Let $\sigma^*$ be a mixed strategy Nash equilibrium. For each player~$i$, every pure
strategy in $\supp(\sigma_i^*)$ yields the same expected payoff against $\sigma_{-i}^*$,
and this common payoff is at least as large as the expected payoff from any pure
strategy outside the support. Formally:
\begin{enumerate}[nosep,label=(\alph*)]
    \item If $s_i, s_i' \in \supp(\sigma_i^*)$, then
    $u_i(s_i, \sigma_{-i}^*) = u_i(s_i', \sigma_{-i}^*)$.
    \item If $s_i \in \supp(\sigma_i^*)$ and $t_i \notin \supp(\sigma_i^*)$, then
    $u_i(s_i, \sigma_{-i}^*) \geq u_i(t_i, \sigma_{-i}^*)$.
\end{enumerate}
\end{theorem}

\begin{proof}
By multilinearity (Remark~\ref{rem:multilinear}),
\[
    u_i(\sigma_i^*, \sigma_{-i}^*) = \sum_{s_i \in S_i} \sigma_i^*(s_i)\, u_i(s_i, \sigma_{-i}^*).
\]
Let $v_i^* = u_i(\sigma_i^*, \sigma_{-i}^*)$ be the equilibrium payoff for player~$i$.

\textbf{Part (b):} Suppose $t_i \in S_i$ satisfies $u_i(t_i, \sigma_{-i}^*) > v_i^*$.
Then the pure strategy $t_i$ (viewed as a degenerate mixed strategy) gives payoff
$u_i(t_i, \sigma_{-i}^*) > v_i^* = u_i(\sigma_i^*, \sigma_{-i}^*)$, contradicting the
NE condition. Hence $u_i(s_i, \sigma_{-i}^*) \leq v_i^*$ for all $s_i$ with equality
impossible for all $s_i$ unless we also have a lower bound. Since
$v_i^*$ is a convex combination of the payoffs $u_i(s_i, \sigma_{-i}^*)$ with
weights $\sigma_i^*(s_i) \geq 0$ summing to~1, if any payoff in the support were
strictly less than $v_i^*$, then some payoff (in or outside the support) must be
strictly greater than $v_i^*$ (to make the weighted average equal $v_i^*$), which
we just showed is impossible. Therefore:
\begin{itemize}[nosep]
    \item For all $s_i \in \supp(\sigma_i^*)$: $u_i(s_i, \sigma_{-i}^*) = v_i^*$.
    \item For all $s_i \notin \supp(\sigma_i^*)$: $u_i(s_i, \sigma_{-i}^*) \leq v_i^*$.
\end{itemize}

\textbf{Part (a):} Immediate from the first bullet point.
\end{proof}

\begin{example}[Mixed NE in Battle of the Sexes]
\label{ex:mixed_bos}
In the Battle of the Sexes (Example~\ref{ex:bos}), let player~1 play $O$ with
probability $p$ and player~2 play $O$ with probability $q$. By the indifference
principle, in a fully mixed NE, player~1 must be indifferent:
\[
    u_1(O, q) = u_1(F, q) \implies 3q + 0(1-q) = 0 \cdot q + 1 \cdot (1-q)
    \implies 3q = 1 - q \implies q = \tfrac{1}{4}.
\]
Similarly, player~2 must be indifferent:
\[
    u_2(p, O) = u_2(p, F) \implies 1 \cdot p + 0(1-p) = 0 \cdot p + 3(1-p)
    \implies p = 3 - 3p \implies p = \tfrac{3}{4}.
\]
The mixed NE is $\sigma_1 = (\tfrac{3}{4}, \tfrac{1}{4})$,
$\sigma_2 = (\tfrac{1}{4}, \tfrac{3}{4})$, with expected payoffs
$(3 \cdot \tfrac{1}{4}, 1 \cdot \tfrac{3}{4}) = (\tfrac{3}{4}, \tfrac{3}{4})$.
Both players fare worse in the mixed NE than in either pure NE.
\end{example}

\begin{rlconnection}[Mixed Strategies and Stochastic Policies]
In RL, a stochastic policy $\pi(a|s)$ assigns probabilities to actions given a state.
This is precisely a mixed strategy in the stage game at state~$s$. The indifference
principle has a direct analog: in an optimal stochastic policy, every action in the
support must yield the same Q-value. If some action had a strictly higher Q-value,
the agent should place all probability on it. This is why $\epsilon$-greedy exploration
is suboptimal --- it assigns probability to actions that may not be in the support of
any Nash equilibrium.
\end{rlconnection}

%% ============================================================
%% SECTION 5: NASH EXISTENCE THEOREM
%% ============================================================

\section{Nash's Existence Theorem}
\label{sec:nash_existence}

The central result of this lesson: every finite game has at least one Nash equilibrium
in mixed strategies. The proof uses Kakutani's fixed point theorem.

\subsection{Kakutani's Fixed Point Theorem}
\label{subsec:kakutani}

\begin{definition}[Upper Hemicontinuity]
\label{def:uhc}
Let $X \subseteq \R^m$ and $Y \subseteq \R^n$. A correspondence (set-valued map)
$F : X \rightrightarrows Y$ is \emph{upper hemicontinuous} (uhc) at $x_0 \in X$ if
for every open set $V \supseteq F(x_0)$, there exists an open set $U \ni x_0$ such
that $F(x) \subseteq V$ for all $x \in U \cap X$.

Equivalently (when $Y$ is compact), $F$ is uhc at $x_0$ if whenever $x_k \to x_0$
and $y_k \in F(x_k)$ with $y_k \to y_0$, we have $y_0 \in F(x_0)$. This is the
\emph{closed graph} characterization.
\end{definition}

\begin{theorem}[Kakutani's Fixed Point Theorem]
\label{thm:kakutani}
Let $K \subseteq \R^n$ be non-empty, compact, and convex. Let $F : K \rightrightarrows K$
be a correspondence such that:
\begin{enumerate}[nosep,label=(\roman*)]
    \item $F(x)$ is non-empty for every $x \in K$.
    \item $F(x)$ is convex for every $x \in K$.
    \item $F$ is upper hemicontinuous.
\end{enumerate}
Then $F$ has a fixed point: there exists $x^* \in K$ with $x^* \in F(x^*)$.
\end{theorem}

\begin{proof}
We provide the standard proof via Brouwer's fixed point theorem.

For each $k \in \N$, let $T_k$ be a triangulation of $K$ with mesh size at most
$1/k$ (the maximum diameter of any simplex in the triangulation). For each vertex $v$
of $T_k$, choose $f_k(v) \in F(v)$ (possible since $F(v)$ is non-empty). Extend
$f_k$ to all of $K$ by convex interpolation on each simplex: if
$x = \sum_{j} \lambda_j v_j$ is the barycentric representation of $x$ in its simplex,
set $f_k(x) = \sum_j \lambda_j f_k(v_j)$.

We claim each $f_k : K \to K$ is continuous. Continuity within each simplex follows
from the linearity of the interpolation. On shared faces of adjacent simplices, the
interpolation agrees because barycentric coordinates are unique on the shared face.
Since $K$ is convex and each $f_k(x)$ is a convex combination of points in $K$,
we have $f_k(x) \in K$.

By Brouwer's fixed point theorem (applicable since $K$ is compact and convex, and
$f_k$ is continuous), each $f_k$ has a fixed point $x_k \in K$ with $f_k(x_k) = x_k$.

Since $K$ is compact, the sequence $(x_k)$ has a convergent subsequence
$x_{k_l} \to x^* \in K$.

We show $x^* \in F(x^*)$. For each $l$, let $\sigma_l$ be the simplex of $T_{k_l}$
containing $x_{k_l}$, with vertices $v_1^l, \ldots, v_p^l$ (where $p \leq n+1$).
Then
\[
    x_{k_l} = f_{k_l}(x_{k_l}) = \sum_{j=1}^p \lambda_j^l f_{k_l}(v_j^l),
    \quad \text{where } x_{k_l} = \sum_{j=1}^p \lambda_j^l v_j^l.
\]
Since the mesh size tends to~0, $\|v_j^l - x_{k_l}\| \leq 1/k_l \to 0$, so
$v_j^l \to x^*$ for each~$j$.

Now, $f_{k_l}(v_j^l) \in F(v_j^l)$. Since $K$ is compact, by passing to a further
subsequence (using a diagonal argument over the finitely many vertices), we may assume
$f_{k_l}(v_j^l) \to y_j$ for some $y_j \in K$.

By the closed-graph characterization of upper hemicontinuity (Definition~\ref{def:uhc}):
$v_j^l \to x^*$ and $f_{k_l}(v_j^l) \to y_j$ with $f_{k_l}(v_j^l) \in F(v_j^l)$
implies $y_j \in F(x^*)$.

The barycentric coordinates $\lambda_j^l \in [0,1]$ lie in a compact set, so
(passing to a further subsequence) $\lambda_j^l \to \lambda_j^*$ with
$\sum_j \lambda_j^* = 1$ and $\lambda_j^* \geq 0$.

Taking the limit in $x_{k_l} = \sum_j \lambda_j^l f_{k_l}(v_j^l)$:
\[
    x^* = \sum_j \lambda_j^* y_j.
\]
Since each $y_j \in F(x^*)$ and $F(x^*)$ is convex, the convex combination
$x^* = \sum_j \lambda_j^* y_j \in F(x^*)$.
\end{proof}

\subsection{Nash's Theorem}
\label{subsec:nash_theorem}

\begin{keyresult}[Nash's Existence Theorem]
\begin{theorem}[Nash, 1950]
\label{thm:nash_existence}
Every finite normal-form game has at least one mixed strategy Nash equilibrium.
\end{theorem}
\end{keyresult}

\begin{proof}
Let $\Gamma = (N, (S_i)_{i \in N}, (u_i)_{i \in N})$ be a finite game with
$|S_i| = m_i$. Define the set of mixed strategy profiles
$\Delta = \prod_{i \in N} \Delta(S_i)$, which is a product of simplices and hence
compact and convex in $\R^{m_1 + \cdots + m_n}$.

\textbf{Step 1: Define the best response correspondence.}
For each player~$i$, the \emph{mixed best response correspondence}
$\BR_i : \Delta_{-i} \rightrightarrows \Delta(S_i)$ is
\[
    \BR_i(\sigma_{-i}) = \argmax_{\sigma_i \in \Delta(S_i)} u_i(\sigma_i, \sigma_{-i}).
\]
Define the joint best response correspondence $\BR : \Delta \rightrightarrows \Delta$ by
\[
    \BR(\sigma) = \BR_1(\sigma_{-1}) \times \BR_2(\sigma_{-2}) \times \cdots
    \times \BR_n(\sigma_{-n}).
\]

\textbf{Step 2: $\BR(\sigma)$ is non-empty.}
For each~$i$, the function $\sigma_i \mapsto u_i(\sigma_i, \sigma_{-i})$ is
continuous (in fact, linear) on the compact set $\Delta(S_i)$. By the extreme value
theorem, the maximum is attained, so $\BR_i(\sigma_{-i}) \neq \emptyset$.

\textbf{Step 3: $\BR(\sigma)$ is convex.}
Fix~$i$ and $\sigma_{-i}$. Let $\sigma_i, \sigma_i' \in \BR_i(\sigma_{-i})$ and
$\lambda \in [0,1]$. Set $\hat{\sigma}_i = \lambda \sigma_i + (1-\lambda)\sigma_i'$.
By linearity of $u_i$ in $\sigma_i$ (Remark~\ref{rem:multilinear}):
\[
    u_i(\hat{\sigma}_i, \sigma_{-i})
    = \lambda\, u_i(\sigma_i, \sigma_{-i}) + (1-\lambda)\, u_i(\sigma_i', \sigma_{-i})
    = \lambda\, v_i^* + (1-\lambda)\, v_i^* = v_i^*,
\]
where $v_i^* = \max_{\tau_i} u_i(\tau_i, \sigma_{-i})$. So $\hat{\sigma}_i$ also
achieves the maximum, hence $\hat{\sigma}_i \in \BR_i(\sigma_{-i})$. Thus
$\BR_i(\sigma_{-i})$ is convex. Since $\BR(\sigma)$ is a product of convex sets,
it is convex.

\textbf{Step 4: $\BR$ is upper hemicontinuous.}
We use the closed-graph characterization. Let $\sigma^k \to \sigma^*$ in $\Delta$
and $\tau^k \in \BR(\sigma^k)$ with $\tau^k \to \tau^*$. We must show
$\tau^* \in \BR(\sigma^*)$, i.e., $\tau_i^* \in \BR_i(\sigma_{-i}^*)$ for each~$i$.

Fix player~$i$. Since $\tau_i^k \in \BR_i(\sigma_{-i}^k)$, for every
$\sigma_i \in \Delta(S_i)$:
\[
    u_i(\tau_i^k, \sigma_{-i}^k) \geq u_i(\sigma_i, \sigma_{-i}^k).
\]
Since $u_i$ is continuous (it is multilinear on the compact set $\Delta$), taking
$k \to \infty$:
\[
    u_i(\tau_i^*, \sigma_{-i}^*) \geq u_i(\sigma_i, \sigma_{-i}^*).
\]
Since this holds for all $\sigma_i$, we have $\tau_i^* \in \BR_i(\sigma_{-i}^*)$.

\textbf{Step 5: Apply Kakutani.}
The correspondence $\BR : \Delta \rightrightarrows \Delta$ satisfies all hypotheses of
Kakutani's theorem (Theorem~\ref{thm:kakutani}): $\Delta$ is non-empty, compact, and
convex; $\BR(\sigma)$ is non-empty and convex for each $\sigma$; and $\BR$ is upper
hemicontinuous. Therefore, there exists $\sigma^* \in \Delta$ with
$\sigma^* \in \BR(\sigma^*)$.

By definition of $\BR$, $\sigma_i^* \in \BR_i(\sigma_{-i}^*)$ for all~$i$,
which by Proposition~\ref{prop:ne_br} (extended to mixed strategies) means $\sigma^*$
is a Nash equilibrium.
\end{proof}

\begin{remark}[Alternative Proof via Brouwer]
\label{rem:brouwer_proof}
An alternative proof avoids Kakutani by constructing a continuous function whose
fixed points are NE. For each player~$i$ and pure strategy $s_i^j$, define the
``excess payoff'':
\[
    g_i^j(\sigma) = \max\big\{0,\; u_i(s_i^j, \sigma_{-i}) - u_i(\sigma_i, \sigma_{-i})\big\}.
\]
Then define $f : \Delta \to \Delta$ by
\[
    f_i^j(\sigma) = \frac{\sigma_i(s_i^j) + g_i^j(\sigma)}
    {1 + \sum_{k=1}^{m_i} g_i^k(\sigma)}.
\]
One verifies: (i) $f$ maps $\Delta$ to $\Delta$; (ii) $f$ is continuous; (iii)
$\sigma^*$ is a fixed point of $f$ if and only if $g_i^j(\sigma^*) = 0$ for all
$i, j$, which holds if and only if $\sigma^*$ is a NE. Brouwer's theorem gives a
fixed point, hence a NE.
\end{remark}

\begin{rlconnection}[Existence Guarantees in MARL]
Nash's theorem guarantees that every finite MARL problem (modeled as a stochastic game
with finite state and action spaces) has at least one Nash equilibrium --- at least in
the stage games. This justifies the goal of Nash-Q learning (Hu \& Wellman, 2003),
which attempts to compute Nash equilibria at each state. However, the theorem is
non-constructive: it guarantees existence but says nothing about efficient computation.
Indeed, computing a Nash equilibrium is PPAD-complete (Daskalakis, Goldberg, Papadimitriou,
2009), meaning no polynomial-time algorithm is known.
\end{rlconnection}

%% ============================================================
%% SECTION 6: MINIMAX THEOREM
%% ============================================================

\section{Von Neumann's Minimax Theorem}
\label{sec:minimax}

\subsection{Two-Player Zero-Sum Games}
\label{subsec:zero_sum}

\begin{definition}[Two-Player Zero-Sum Game]
\label{def:zero_sum}
A two-player game $(N, (S_1, S_2), (u_1, u_2))$ is \emph{zero-sum} if
$u_1(s) + u_2(s) = 0$ for all $s \in S$. It is specified by a single payoff matrix
$A \in \R^{m \times n}$ where $u_1(i,j) = A_{ij}$ and $u_2(i,j) = -A_{ij}$.
Player~1 (the ``row player'') maximizes; player~2 (the ``column player'') minimizes.
\end{definition}

\begin{definition}[Maximin and Minimax Values]
\label{def:maximin}
For a zero-sum game with payoff matrix $A$, define:
\begin{align*}
    \underline{v} &= \max_{x \in \Delta_m} \min_{y \in \Delta_n} x^\top A y
    \quad \text{(maximin value: player~1's guarantee)}, \\
    \overline{v} &= \min_{y \in \Delta_n} \max_{x \in \Delta_m} x^\top A y
    \quad \text{(minimax value: player~2's guarantee)},
\end{align*}
where $\Delta_k$ denotes the $(k-1)$-simplex in $\R^k$.
\end{definition}

\begin{lemma}[Maximin $\leq$ Minimax]
\label{lem:maximin_leq}
For any matrix $A \in \R^{m \times n}$, $\underline{v} \leq \overline{v}$.
\end{lemma}

\begin{proof}
For any fixed $x \in \Delta_m$ and $y \in \Delta_n$:
\[
    \min_{y' \in \Delta_n} x^\top A y' \leq x^\top A y
    \leq \max_{x' \in \Delta_m} (x')^\top A y.
\]
Taking $\max$ over $x$ on the left and $\min$ over $y$ on the right:
\[
    \underline{v} = \max_x \min_{y'} x^\top A y'
    \leq \min_y \max_{x'} (x')^\top A y = \overline{v}. \qedhere
\]
\end{proof}

\begin{keyresult}[Von Neumann's Minimax Theorem]
\begin{theorem}[Von Neumann, 1928]
\label{thm:minimax}
For every matrix $A \in \R^{m \times n}$,
\[
    \max_{x \in \Delta_m} \min_{y \in \Delta_n} x^\top A y
    = \min_{y \in \Delta_n} \max_{x \in \Delta_m} x^\top A y.
\]
The common value $v^* = \underline{v} = \overline{v}$ is called the \emph{value}
of the game.
\end{theorem}
\end{keyresult}

\begin{proof}
We give a proof using LP duality. By Lemma~\ref{lem:maximin_leq}, it suffices to
show $\underline{v} \geq \overline{v}$.

\textbf{Step 1: Reformulate as linear programs.}
Player~1's maximin problem is:
\[
    \underline{v} = \max_{x \in \Delta_m} \min_{y \in \Delta_n} x^\top A y.
\]
Since $y \mapsto x^\top A y$ is linear in $y$ and $\Delta_n$ is a simplex, the minimum
over $y$ is attained at a vertex $e_j$ (the $j$-th standard basis vector), giving
$\min_y x^\top A y = \min_{j=1,\ldots,n} (A^\top x)_j$.

Thus $\underline{v} = \max_{x \in \Delta_m} \min_{j} (A^\top x)_j$. Introducing an
auxiliary variable $v$, this is equivalent to the linear program:
\begin{align}
    \text{(P)} \quad \max\;&\; v \notag \\
    \text{s.t.}\;&\; A^\top x \geq v\, \mathbf{1}, \quad x \in \Delta_m, \quad v \in \R,
    \label{eq:primal}
\end{align}
where $\mathbf{1} = (1, \ldots, 1)^\top \in \R^n$ and the inequality is componentwise.

\textbf{Step 2: The dual LP.}
Similarly, player~2's minimax problem gives:
\[
    \overline{v} = \min_{y \in \Delta_n} \max_{i} (Ay)_i.
\]
This is equivalent to:
\begin{align}
    \text{(D)} \quad \min\;&\; w \notag \\
    \text{s.t.}\;&\; Ay \leq w\, \mathbf{1}, \quad y \in \Delta_n, \quad w \in \R.
    \label{eq:dual}
\end{align}

\textbf{Step 3: LP duality.}
We show that (P) and (D) are LP duals of each other. Rewrite (P) in standard form.
Let $v = \underline{v}$. The constraints of (P) are:
$A^\top x - v\, \mathbf{1} \geq 0$, $\mathbf{1}^\top x = 1$, $x \geq 0$.
This is a linear program in the variables $(x, v) \in \R^{m+1}$.

The dual of this LP (after careful bookkeeping of the dual variables associated
with each primal constraint) yields precisely the LP (D). By strong LP duality
(both problems are feasible --- take $x = (1/m, \ldots, 1/m)^\top$ and
$y = (1/n, \ldots, 1/n)^\top$ --- and bounded), the optimal values are equal:
\[
    \underline{v} = \text{opt}(\text{P}) = \text{opt}(\text{D}) = \overline{v}. \qedhere
\]
\end{proof}

\begin{corollary}[NE in Zero-Sum Games]
\label{cor:ne_zero_sum}
In a two-player zero-sum game with payoff matrix $A$ and value $v^*$:
\begin{enumerate}[nosep,label=(\alph*)]
    \item $(x^*, y^*)$ is a Nash equilibrium if and only if $x^*$ is a maximin
    strategy and $y^*$ is a minimax strategy.
    \item All NE yield the same payoff: $u_1 = v^*$, $u_2 = -v^*$.
    \item NE strategies are interchangeable: if $(x^*, y^*)$ and $(x^{**}, y^{**})$
    are both NE, then so are $(x^*, y^{**})$ and $(x^{**}, y^*)$.
\end{enumerate}
\end{corollary}

\begin{proof}
\textbf{(a)} $(\Rightarrow)$: Let $(x^*, y^*)$ be a NE. Then $x^*$ is a best response
to $y^*$, so $(x^*)^\top A y^* = \max_x x^\top A y^*$. Similarly, $y^*$ minimizes
$x^\top A y$ against $x^*$, so $(x^*)^\top A y^* = \min_y (x^*)^\top A y$. Therefore:
\[
    \max_x x^\top A y^* = (x^*)^\top A y^* = \min_y (x^*)^\top A y.
\]
But $\min_y (x^*)^\top A y \leq \underline{v} \leq \overline{v}
\leq \max_x x^\top A y^*$,
and by the minimax theorem $\underline{v} = \overline{v} = v^*$, so all inequalities
are equalities. Hence $\min_y (x^*)^\top A y = v^*$ (so $x^*$ achieves the maximin)
and $\max_x x^\top A y^* = v^*$ (so $y^*$ achieves the minimax).

$(\Leftarrow)$: If $x^*$ achieves the maximin and $y^*$ achieves the minimax, then
\[
    (x^*)^\top A y^* \geq \min_y (x^*)^\top A y = v^*
    = \max_x x^\top A y^* \geq (x^*)^\top A y^*,
\]
so $(x^*)^\top A y^* = v^*$. For any $x$: $x^\top A y^* \leq \max_{x'} (x')^\top A y^*
= v^* = (x^*)^\top A y^*$, so $x^*$ is a best response to $y^*$. Similarly, $y^*$
is a best response to $x^*$. Hence $(x^*, y^*)$ is a NE.

\textbf{(b)} Follows from part (a): every NE strategy pair achieves the value $v^*$.

\textbf{(c)} If $x^*$ and $x^{**}$ are both maximin and $y^*$, $y^{**}$ both minimax,
then by part (a), any combination of a maximin and a minimax strategy forms a NE.
\end{proof}

\begin{rlconnection}[Minimax in RL]
The minimax theorem is the theoretical foundation of:
\begin{itemize}[nosep]
    \item \textbf{Minimax-Q learning} (Littman, 1994): at each state of a two-player
    zero-sum stochastic game, the Bellman operator computes
    $V(s) = \max_x \min_y x^\top Q(s, \cdot, \cdot) y$, which is well-defined
    by the minimax theorem.
    \item \textbf{AlphaGo/AlphaZero}: the MCTS search implicitly approximates
    minimax values in the game tree.
    \item \textbf{GANs}: the generator-discriminator game is a continuous zero-sum
    game $\min_G \max_D V(D,G)$. The minimax theorem (extended to continuous
    strategy spaces) justifies seeking saddle points.
\end{itemize}
\end{rlconnection}

\subsection{Connection to Linear Programming}
\label{subsec:lp}

\begin{proposition}[Solving Zero-Sum Games via LP]
\label{prop:lp_zerosum}
The maximin strategy $x^*$ and the game value $v^*$ can be computed by solving the
LP~\eqref{eq:primal}. The minimax strategy $y^*$ is recovered from the dual
LP~\eqref{eq:dual}. Both LPs have polynomial-time algorithms (e.g., interior point
methods), so computing NE in two-player zero-sum games is in~P.
\end{proposition}

\begin{proof}
This follows directly from the reformulation in the proof of
Theorem~\ref{thm:minimax}. The LP~\eqref{eq:primal} has $m + 1$ variables and
$n + m + 1$ constraints (including non-negativity and the simplex constraint), so
it can be solved efficiently.
\end{proof}

\begin{warning}[Zero-Sum vs.\ General-Sum Complexity]
While NE in two-player zero-sum games can be found in polynomial time via LP,
computing a Nash equilibrium in a general two-player game is PPAD-complete
(Chen, Deng, Teng, 2009; Daskalakis, Goldberg, Papadimitriou, 2009). This
complexity-theoretic barrier is a fundamental challenge for MARL: even
\emph{finding} the solution concept that algorithms like Nash-Q target is
computationally intractable in general.
\end{warning}

%% ============================================================
%% SECTION 7: CORRELATED EQUILIBRIUM
%% ============================================================

\section{Correlated Equilibrium}
\label{sec:correlated}

\begin{definition}[Correlated Equilibrium]
\label{def:ce}
A \emph{correlated equilibrium} (CE) of a finite game $\Gamma = (N, (S_i), (u_i))$
is a probability distribution $\pi$ over $S = \prod_{i \in N} S_i$ such that for
every player~$i$ and every pair of strategies $s_i, s_i' \in S_i$:
\[
    \sum_{s_{-i} \in S_{-i}} \pi(s_i, s_{-i})\, u_i(s_i, s_{-i})
    \geq \sum_{s_{-i} \in S_{-i}} \pi(s_i, s_{-i})\, u_i(s_i', s_{-i}).
\]
Equivalently, if a mediator draws $s = (s_1, \ldots, s_n)$ from $\pi$ and
privately recommends $s_i$ to player~$i$, no player has an incentive to deviate
from the recommendation (assuming the others follow theirs).
\end{definition}

\begin{theorem}[Every Nash Equilibrium Is a Correlated Equilibrium]
\label{thm:ne_is_ce}
If $\sigma^* = (\sigma_1^*, \ldots, \sigma_n^*)$ is a mixed strategy NE, then the
product distribution $\pi(s) = \prod_{i \in N} \sigma_i^*(s_i)$ is a correlated
equilibrium.
\end{theorem}

\begin{proof}
Fix player~$i$ and strategies $s_i, s_i' \in S_i$. The CE incentive constraint
requires:
\[
    \sum_{s_{-i}} \pi(s_i, s_{-i})\, u_i(s_i, s_{-i})
    \geq \sum_{s_{-i}} \pi(s_i, s_{-i})\, u_i(s_i', s_{-i}).
\]
Since $\pi$ is a product distribution, $\pi(s_i, s_{-i}) = \sigma_i^*(s_i) \cdot
\prod_{j \neq i} \sigma_j^*(s_j)$. If $\sigma_i^*(s_i) = 0$, both sides are~0
and the inequality holds. If $\sigma_i^*(s_i) > 0$, dividing both sides by
$\sigma_i^*(s_i)$ gives:
\[
    \sum_{s_{-i}} \bigg(\prod_{j \neq i} \sigma_j^*(s_j)\bigg) u_i(s_i, s_{-i})
    \geq \sum_{s_{-i}} \bigg(\prod_{j \neq i} \sigma_j^*(s_j)\bigg) u_i(s_i', s_{-i}),
\]
i.e., $u_i(s_i, \sigma_{-i}^*) \geq u_i(s_i', \sigma_{-i}^*)$.
Since $s_i \in \supp(\sigma_i^*)$, the indifference principle
(Theorem~\ref{thm:indifference}) gives
$u_i(s_i, \sigma_{-i}^*) \geq u_i(s_i', \sigma_{-i}^*)$ for all $s_i'$, which is
exactly what we need.
\end{proof}

\begin{proposition}[CE as a Linear Program]
\label{prop:ce_lp}
The set of correlated equilibria of a finite game is a convex polytope, defined by the
linear incentive-compatibility constraints in Definition~\ref{def:ce} together with
$\pi(s) \geq 0$ for all $s$ and $\sum_s \pi(s) = 1$. In particular, optimizing any
linear objective over the CE set is a linear program.
\end{proposition}

\begin{proof}
The CE constraints are:
\[
    \sum_{s_{-i}} \pi(s_i, s_{-i})\big[u_i(s_i, s_{-i}) - u_i(s_i', s_{-i})\big] \geq 0
    \quad \text{for all } i, s_i, s_i'.
\]
These are linear in the variables $\{\pi(s)\}_{s \in S}$. Combined with the simplex
constraints $\pi(s) \geq 0$ and $\sum_s \pi(s) = 1$, the feasible set is a polytope
(intersection of finitely many half-spaces). Any linear function of $\pi$ can be
optimized over this polytope by solving an LP.
\end{proof}

\begin{example}[Correlated Equilibrium in Battle of the Sexes]
\label{ex:ce_bos}
In the Battle of the Sexes (Example~\ref{ex:bos}), consider the ``traffic light''
correlated equilibrium: a mediator recommends $(O,O)$ with probability $1/2$ and
$(F,F)$ with probability $1/2$. Under this distribution $\pi$:
\[
    \pi(O,O) = 1/2, \quad \pi(O,F) = 0, \quad \pi(F,O) = 0, \quad \pi(F,F) = 1/2.
\]
Check the CE condition for player~1, given recommendation $O$ (so $s_1 = O$):
\[
    \pi(O,O)\, u_1(O,O) = \tfrac{1}{2} \cdot 3 = \tfrac{3}{2}
    \geq \pi(O,O)\, u_1(F,O) = \tfrac{1}{2} \cdot 0 = 0. \quad \checkmark
\]
Given recommendation $F$:
\[
    \pi(F,F)\, u_1(F,F) = \tfrac{1}{2} \cdot 1 = \tfrac{1}{2}
    \geq \pi(F,F)\, u_1(O,F) = \tfrac{1}{2} \cdot 0 = 0. \quad \checkmark
\]
Similarly for player~2. The expected payoffs are
$(1/2)(3) + (1/2)(1) = 2$ for player~1 and $(1/2)(1) + (1/2)(3) = 2$ for player~2,
which Pareto-dominates the mixed NE payoffs of $(3/4, 3/4)$.
\end{example}

\begin{rlconnection}[Correlated Equilibrium and MARL]
Correlated equilibrium has several advantages over Nash equilibrium for MARL:
\begin{itemize}[nosep]
    \item \textbf{Computational tractability}: finding a CE is an LP, polynomial in
    the game size. Finding a NE is PPAD-complete.
    \item \textbf{Natural emergence}: simple no-regret learning dynamics (e.g., regret
    matching) converge to the set of CE, not NE. This makes CE a more natural
    prediction for what independent learners will achieve.
    \item \textbf{Better welfare}: CE can achieve higher social welfare than any NE,
    as the Battle of the Sexes example illustrates.
\end{itemize}
\end{rlconnection}

%% ============================================================
%% SECTION 8: EXTENSIVE-FORM GAMES
%% ============================================================


\end{document}